\documentclass[11pt]{amsart}

\usepackage{amsmath}
\usepackage{amssymb}
\usepackage{amsthm}
%\usepackage{psfig}

\begin{document}
\baselineskip=15pt
\textheight=8.4in
\parindent=0pt
\def\sk {\hskip .5cm}
\def\skv {\vskip .12cm}
\def\cos {\mbox{cos}}
\def\sin {\mbox{sin}}
\def\tan {\mbox{tan}}
\def\intl{\int\limits}
\def\lm{\lim\limits}
\newcommand{\frc}{\displaystyle\frac}
\def\xbf{{\mathbf x}}
\def\fbf{{\mathbf f}}
\def\gbf{{\mathbf g}}

\def\Ker{{\rm Ker\,}}
\def\phi{\varphi}

\def\dbA{{\mathbb A}}
\def\dbB{{\mathbb B}}
\def\dbC{{\mathbb C}}
\def\dbD{{\mathbb D}}
\def\dbE{{\mathbb E}}
\def\dbF{{\mathbb F}}
\def\dbG{{\mathbb G}}
\def\dbH{{\mathbb H}}
\def\dbI{{\mathbb I}}
\def\dbJ{{\mathbb J}}
\def\dbK{{\mathbb K}}
\def\dbL{{\mathbb L}}
\def\dbM{{\mathbb M}}
\def\dbN{{\mathbb N}}
\def\dbO{{\mathbb O}}
\def\dbP{{\mathbb P}}
\def\dbQ{{\mathbb Q}}
\def\dbR{{\mathbb R}}
\def\dbS{{\mathbb S}}
\def\dbT{{\mathbb T}}
\def\dbU{{\mathbb U}}
\def\dbV{{\mathbb V}}
\def\dbW{{\mathbb W}}
\def\dbX{{\mathbb X}}
\def\dbY{{\mathbb Y}}
\def\dbZ{{\mathbb Z}}


\def\char{{\rm char}}
\def\wt{{\rm wt}}
\def\Aut{{\rm Aut}}
\def\Hom{{\rm Hom}}
\def\End{{\rm End}}
\def\Im{{\rm Im}}
\def\GL{{\rm GL}}
\def\SL{{\rm SL}}
\def\SNF{{\rm SNF}}
\def\tr {{\rm tr}}
\def\rk{{\rm rk}}
\def\lam{{\lambda}}

\def\la{{\langle}}
\def\ra{{\rangle}}

\bf\centerline{Homework \# 6, to be submitted on Canvas by 6pm on Sat, Mar 14th.}\rm
\vskip .2cm
{\bf Plan for the next 2 classes (Mar 10 and 12):} Character tables (continued) and orthogonality relations for characters
(see Lectures 16-19 from Math 5657 notes).
Decompositions of the regular representations as a direct sum of irreducibles, time permitting (see Lectures 20-21 from Math 5657 notes).
\vskip .1cm

%{\bf Note on hints:} Some hints are given at the end of the assignment, each on a separate page.
%Problems (or parts of problems) for which a hint is available at the end are marked with *.

\skv
{\bf Problem 1:} Let $V$ be an inner product space.
\begin{itemize}
\item[(a)] (practice, no need to turn in) Prove that for any subspace $W$ of $V$ we have $V=W\oplus W^{\perp}$.
\item[(b)] Now assume that $(\rho,V)$ is a unitary representation of a group $G$. Prove that if $W\subseteq V$
is a subrepresentation, then so is $W^{\perp}$.
\end{itemize}
\skv
{\bf Problem 2:} Let $A=(a_{ij})\in Mat_n(\dbC)$. The matrix  $A^*=\overline{A^T}$ (the conjugate transpose of $A$) is called
the {\it adjoint matrix} of $A$. Note that $(A^*)_{ij}=\overline{a_{ji}}$. The matrix $A$ is called unitary if $A^*A=AA^*=I$, that is, if $A$ is invertible and $A^*=A^{-1}$.
\begin{itemize}
\item[(1)] Prove that the following are equivalent for a given $A$:
\begin{itemize}
\item[(a)] $A$ is unitary;
\item[(b)] Rows of $A$ form an orthonormal basis of $\dbC^n$ (with respect to the Hermitian dot product);
\item[(c)] Columns of $A$ form an orthonormal basis of $\dbC^n$ (with respect to the Hermitian dot product).
\end{itemize}
\item[(2)] Now assume that $(\rho,V)$ is a unitary representation of a group $G$, and let $\beta$ be an orthonormal basis for $V$.
Prove that $(\rho,V)$ is unitary $\iff$ $[\rho(g)]_{\beta}$ is unitary for each $g\in G$.
\end{itemize}
\skv

{\bf Problem 3:} The goal of this problem is to prove the general case of Maschke's Theorem: \it If $G$ is a finite group and $F$ is any field with
$\char(F)\nmid |G|$, then any representation of $G$ over $F$ is completely reducible. \rm

\sk By Theorem~11.1 from class, we just need to show that every representation of $G$ over $F$ has the complement property:

{\bf Lemma:} \it Let $G$ and $F$ be as above, $(\rho,V)$ a representation of $G$ over $F$ and $W$ a $G$-invariant subspace. Then there exists a $G$-invariant subspace $U$ such that $V=W\oplus U$. \rm

\sk To prove the lemma, consider the following maps. Choose any subspace $Z$ of $V$ such that $V=W\oplus Z$ and
let $P: V\to W$ be the projection onto $W$ along $Z$, that is, $P$ is the unique linear map such that $P(w)=w$ for all $w\in W$ and $P(Z)=0$. Now define $Q:V\to V$ by
$$Q(v)=\frac{1}{|G|}\sum\limits_{g\in G}\rho(g)^{-1} P \rho(g) (v).$$
Prove that 
\begin{itemize}
\item[(i)] $Q(w)=w$ for all $w\in W$ and $\Im(Q)=W$. Deduce that $Q^2=Q$.
\item[(ii)] $\Ker(Q)$ is $G$-invariant
\end{itemize}
Deduce from (i) that $V=W\oplus \Ker(Q)$ and therefore $U=\Ker(Q)$ has the desired property.
\skv
{\bf Problem 4:}
\begin{itemize}
\item[(a)] Let $M$ and $Q$ be modules over some ring $R$. Suppose that $M$ is completely reducible and $Q$ is a quotient of $M$.
Prove that $M$ has a submodule isomorphic to $Q$.
\item[(b)] Use (a) and Maschke's theorem to prove thatif $G$ is a finite group and $F$ is a field with $\char(F)\nmid |G|$,
then for any irreducible representation $V$ of $G$, there is a subrepresentation of the regular representation $F[G]$ equivalent to $V$.

Recall that the regular representation is the permutation representation associated to the left multiplication action of $G$ on itself.
As an $F[G]$-module, this representation is just a free module of rank~$1$.
\end{itemize}
\skv
{\bf Problem 5:} In Lecture~12 we proved that if $G$ is an abelian group, then any finite-dimensional irreducible representation of $G$ over an algebraically closed field is one-dimensional. 
\begin{itemize}
\item[(a)] Let $n>2$ be an integer. Construct an finite-dimensional irreducible representation of $\dbZ/n\dbZ$ over $\dbR$ (reals) which is not one-dimensional. {\bf Hint:} Recall that we completely described representations of cyclic groups over arbitrary fields in Lecture 10.
\item[(b)] Prove that any irreducible representation of $\dbZ/2\dbZ$ over any field is one-dimensional.
\item[(c)] Describe (with proof) all finite abelian groups $G$ such that any irreducible representation of $G$ over any field is one-dimensional.
\end{itemize}
{\bf Problem 6:} In all parts of this problem include ALL relevant computations.

Let $G=D_{2n}$ be the dihedral group of order $2n$. Recall that $G$ has a presentation 
$G=\la r,s \mid r^n=1, s^2=1, srs^{-1}=r^{-1}\ra$. 
\begin{itemize}
\item[(a)] Let $\omega\in\dbC$ be an $n^{\rm th}$ root of unity (that is, $\omega^n=1$). Prove that $G$ has a representation 
$\rho_{\omega}:G\to \GL_2(\dbC)$ such that $\rho(r)=\begin{pmatrix}\omega & 0\\ 0 &\omega^{-1}
\end{pmatrix}$ and $\rho(s)=\begin{pmatrix}0 & 1\\ 1 &0\end{pmatrix}$
\item[(b)] Prove that the representation $\rho_{\omega}$ in part (a) is irreducible if and only if $\omega\neq \pm 1$.
\item[(c)] Let $\omega_1$ and $\omega_2$ be $n^{\rm th}$ roots of unity different from $\pm 1$. Prove that 
$\rho_{\omega_1}\cong \rho_{\omega_2}$ (as representations of $G$) $\iff$ $\omega_2=\omega_1$ or $\omega_2=\omega_1^{-1}$.
\item[(d)] Since we can think of isometries of a regular $2n$-gon as invertible linear operators on $\dbR^2$, we get a 2-dimensional representation of $G$ ``for free'' (simply representing the elements of $D_{2n}$ by their matrices with respect to the standard basis).
If we assume that $r$ is the counterclockwise rotation by $\frac{2\pi}{n}$ and $s$ is the reflection with respect to the $x$-axis, 
the corresponding $\rho$ is given by $\rho(r)=\begin{pmatrix}\cos\frac{2\pi}{n} & -\sin\frac{2\pi}{n}\\ \sin\frac{2\pi}{n} &\cos\frac{2\pi}{n}\end{pmatrix}$ and $\rho(s)=\begin{pmatrix}1 & 0\\ 0 &-1\end{pmatrix}$. Note that technically $\rho$ is a real representation
(a homomorphism from $G$ to $GL_2(\dbR)$), but since every $2\times 2$ real matrix can be thought of as a  $2\times 2$ complex matrix,
we can think of $\rho$ as a complex representation. Prove that this $\rho$ is equivalent to some $\rho_{\omega}$ from part (a),
explicitly find such $\omega$ as well as a matrix $T\in GL_2(\dbC)$ such that $T^{-1}\rho(g) T=\rho_{\omega}(g)$ for all $g\in G$. 
\item[(e)] Recall that the number of one-dimensional complex representations of a finite group $G$ is equal to $|G^{ab}|$. 
Prove that $[G,G]=\la r^2\ra$ (hint: there is no need to compute the commutator of two arbitrary elements). Deduce that $|G^{ab}|=2$
if $n$ is odd and $|G^{ab}|=4$ if $n$ is even. Then describe all one-dimensional complex representations of $G$ explicitly
(it is enough to specify the images of $r$ and $s$ under these representations).  
\end{itemize}
We will soon prove that every complex irreducible representation of $D_{2n}$ is 1-dimensional or 2-dimensional and, in the latter case,
isomorphic to some $\rho_{\omega}$ as in part (a).
\skv
{\bf Problem 7:} Let $(\alpha,V)$ and $(\beta,W)$ be representations of the same group over the same field. The {\it tensor product} of these representations is the representation $(\rho, V\otimes W)$ where $\rho(g)=\alpha(g)\otimes \beta(g)$ for each $g\in G$ (in the notations
from Problem~6 in HW\#5). In other words, $\rho(g)\in GL(V\otimes W)$ is the unique linear map such that $\rho(g)(v\otimes w)=
(\alpha(g)v)\otimes (\beta(g)w)$ for all $v\in V$ and $w\in W$.
\begin{itemize}
\item[(a)] Prove that $(\rho, V\otimes W)$ is indeed a representation, that is, $\rho: G\to GL(V\otimes W)$ is a homomorphism.
\item[(b)] Prove that if $(\alpha,V)$ is NOT irreducible and $W\neq 0$, then $(\rho,V\otimes W)$ is not irreducible either.
\item[(c)] Now prove that if $(\alpha,V)$ is irreducible and $\dim(W)=1$, then $(\rho,V\otimes W)$ is irreducible.
\end{itemize}
\skv

\end{document}
